% !TEX root = ../aa_SoM_Chapters.tex

%% HARRIS: Chapter 10

% ö = \%c3\%b6
% ä = \%C3\%A4

\xdef\sectiontitle{\IIsectiontitle} %aiheuttama

%%%%%%%%%%%%%%%
\begin{frame}[label=2_4_a]%[label=2_3_TOC1]
\frametitle{\texorpdfstring{\culine[\sectiontitleulinecolor]{\sectiontitle}}{\sectiontitle} }

%\vspace{0.5cm}
%\begin{itemize}[leftmargin=0.5cm,itemsep=3mm]
%    \item[2.1] \hyperlink{2_1_a}{\IIaSubsectiontitle}
%    \item[2.2] \hyperlink{2_2_a}{\IIbSubsectiontitle}	
%    \item[2.3] \hyperlink{2_3_a}{\CDAlert[blue]{\IIcSubsectiontitle}}	
%    \item[2.4] \hyperlink{2_4_a}{\IIdSubsectiontitle}
%    \item[2.5] \hyperlink{2_5_a}{\IIeSubsectiontitle}
%    \item[2.6] \hyperlink{2_6_a}{\IIfSubsectiontitle}
%    \item[2.7] \hyperlink{2_7_a}{\IIgSubsectiontitle}
%    \item[2.8] \hyperlink{2_8_a}{\IIhSubsectiontitle}
%    \item[2.9] \hyperlink{2_9_a}{\IIiSubsectiontitle}
%\end{itemize}

\vspace{0.2cm}
\begin{itemize}[leftmargin=0.5cm,itemsep=3mm]
    \item[2.1] \hyperlink{2_1_a}{\IIaaSubsectiontitle}
    \item[2.2] \hyperlink{2_2_a}{\IIaSubsectiontitle}
    \item[2.3] \hyperlink{2_3_a}{\IIbSubsectiontitle}	
    \item[2.4] \hyperlink{2_4_a}{\CDAlert[blue]{\IIcSubsectiontitle}}	
    \item[2.5] \hyperlink{2_5_a}{\IIdSubsectiontitle}
    \item[2.6] \hyperlink{2_6_a}{\IIeSubsectiontitle}
    \item[2.7] \hyperlink{2_7_a}{\IIfSubsectiontitle}
    \item[2.8] \hyperlink{2_8_a}{\IIgSubsectiontitle}
    \item[2.9] \hyperlink{2_9_a}{\IIhSubsectiontitle}
    \item[2.10] \hyperlink{2_10_a}{\IIiSubsectiontitle}
\end{itemize}


\endofslide \end{frame}
%%%%%%%%%%%%%%%

\xdef\sectiontitle{\IIcSubsectiontitle} 

%%%%%%%%%%%%%%%
\begin{frame}
\frametitle{\texorpdfstring{\culine[\sectiontitleulinecolor]{\sectiontitle}}{\sectiontitle} }

\begin{itemize}
	\item Matter \appear{is an aggregate of a \CDAlert[red]{very large number of atoms}}

\item Appears in \CDAlert[black]{three physical} states \appear{or \CDAlert[black]{phases}:} \appear{\CDAlert[violet]{gases}}\appear{, \CDAlert[blue]{liquids}}\appear{, \CDAlert[red]{solids}} \appear{(the $4^{\text {th}}$ one, \CDAlert[fuchsia]{plasma}, is not discussed here)}

\item In gases the average distance between molecules  \appear{is much greater than sizes of the molecules} 
\implies{ \CDAlert[violet]{Molecules retain their individuality} }

\item \CDAlert[red]{In solids, atoms} (or molecules) are \CDAlert[red]{tightly packed}\appear{, and \CDAlert[red]{held at fixed positions} by forces}\appear{, of electromagnetic origin}\appear{, which are of the same order of magnitude as those involved in the molecular bonding processes} 

\item In solids the atoms are not isolated\appear{, but their \CDAlert[red]{properties are modified by nearby atoms}}

\end{itemize}
% two last form condensed matter .. add a note
%Three phases of water (Richard Feynman, 1964 Lectures in Physics, Voll)

\endofslide 
%\endofsection
\end{frame}
%%%%%%%%%%%%%%%


%%%%%%%%%%%%%%%
\begin{frame}
\frametitle{\texorpdfstring{\culine[\sectiontitleulinecolor]{\sectiontitle}}{\sectiontitle} }

\begin{itemize}
\item \CDAlert[blue]{Liquids} fall in between gases and solids 
\item[] \begin{itemize}
	\item In liquids the molecules form \CDAlert[blue]{temporary short-range} bonds \appear{that are continually broken and reformed} \appear{(due to kinetic energy)}
\end{itemize}

\item The modern field of \CDAlert[blue]{condensed matter physics} \appear{\Alert[blue]{concerns the study of \CDAlert[blue]{both solids and liquids}}}

\item As the name suggests\appear{, the field of \CDAlert[red]{solid-state physics}} \appear{\Alert[red]{'only' focuses on understanding solids} and their properties}

\end{itemize}
%

\endofslide 
%\endofsection
\end{frame}
%%%%%%%%%%%%%%%

%%%%%%%%%%%%%%%
\begin{frame}
\frametitle{\texorpdfstring{\culine[\sectiontitleulinecolor]{\sectiontitle}}{\sectiontitle} }
\framesubtitle{Crystallization}

\begin{itemize}
	\item \CDAlert[red]{Most solids are crystalline}\appear{, with atoms, ions, or molecules} \appear{falling into regular}\appear{, repeated three dimensional patterns}\appear{, forming a crystal lattice} 
	\implies{ Solids have \CDAlert[red]{long-range order}}
	
\item \CDAlert[tunired]{Amorphous solids} \appear{(e.g. glasses, rocks and rubber)} \appear{lack the definite arrangements of their member particles}\appear{, i.\,e. constituting atoms/molecules are not periodically arranged }

\item If you cool down a liquid slowly\appear{, the kinetic energy of molecules is reduced and molecules arrange themselves in regular crystalline order}\appear{, maximum number of bonds are formed resulting in minimized potential energy}

\item If you \CDAlert[tunired]{cool the liquid rapidly}\appear{, so that the internal energy is removed before the molecules have the chance to arrange themselves}\appear{, you'll get an \CDAlert[tunired]{amorphous} solid}\appear{, a "snapshot" of a liquid}
\end{itemize}

\endofslide 
%\endofsection
\end{frame}
%%%%%%%%%%%%%%%

%%%%%%%%%%%%%%%
\begin{frame}
\frametitle{\texorpdfstring{\culine[\sectiontitleulinecolor]{\sectiontitle}}{\sectiontitle} }
\framesubtitle{Crystallization}

\begin{itemize}
	\item Most common solids are \CDAlert[red]{polycrystalline}\appear{, i.\,e. they consist of a collection of single crystals}

\item Single-crystal sizes may be in the order of a fraction of a millimeter 

\item The crystal size depends e.g. on the temperature history of the formation process

\item A single crystal may be said to be built based on the smallest building block\appear{, \CDAlert[red]{a unit cell}}\appear{, which is repeated throughout the whole crystal}\appear{, \CDAlert[tunired]{defining its symmetry and regularity}}

\item The type of unit cell \appear{\CDAlert[red]{depends on the type of bonding} \CDAlert[red]{between atoms}, ions or molecules in the crystal}
\end{itemize}

\endofslide 
%\endofsection
\end{frame}
%%%%%%%%%%%%%%%

%%%%%%%%%%%%%%%
\begin{frame}
\frametitle{\texorpdfstring{\culine[\sectiontitleulinecolor]{\sectiontitle}}{\sectiontitle} }
\framesubtitle{Crystal deformations, stress and strain}

%Crystal defects
\begin{itemize}
	\item In a perfect crystal \appear{each atom has a definite equilibrium location} 
	
	\item Actual crystals are not perfect\appear{, but rather exhibit defects:} \appear{\CDAlert[red]{missing atoms}}\appear{, \CDAlert[red]{atoms out of place}}\appear{, \CDAlert[red]{irregularities} in the spacing of rows of atoms}\appear{, and presence of \CDAlert[red]{impurities}}
	
	\item A \CDAlert[red]{crystal under stress} may be deformed \appear{(i.\,e. is \CDAlert[tunired]{strained})} \appear{irreversibly when \CDAlert[red]{dislocations} in its structure shift the atomic positions}
	
\item Real materials \appear{(not ultrapure ones)} \appear{are strained more easily than what the simple theory predicts}\appear{, which is caused by presence of defects}\appear{, such as edge dislocations in the solid}

%Initial configuration of a crystal with an edge dislocation
\item Let's briefly consider how this happens \appear{(leaving the more involved treatment to the course on solid-state physics...)}
%\item The dislocation moves to the right as the atoms in the layer under it successively shift their bonds with those of the upper layer, one line at the time
	
\end{itemize}


\endofslide 
%\endofsection
\end{frame}
%%%%%%%%%%%%%%%


%%%%%%%%%%%%%%%
\begin{frame}
\frametitle{\texorpdfstring{\culine[\sectiontitleulinecolor]{\sectiontitle}}{\sectiontitle} }
\framesubtitle{Crystal defects}
%Crystal defects

\begin{itemize}
	
%Initial configuration of a crystal with an edge dislocation
\item \CDAlert[fuchsia]{Dislocation} is an abrupt change in the arrangement of atoms \appear{(Figure)}%
\appear{\xdef\loc{south}%
\begin{tikzpicture}[remember picture,overlay]% anchor=south
    \node (A) [yshift=-0.05*\figYshift,xshift=-3.5*\figXshift, anchor=\loc] at (current page.\loc) {\includegraphics[page=1,height=5.5cm]{Figures_booklet/SOM_10_Bonding_figures/Tikz_SOM_Bonding_Mechanical_edge_dislocation_fixed}};%Tikz_SOM_Bonding_Mechanical_dislocation_movement
\end{tikzpicture}}%
\item Presence of dislocations \appear{(and other defects)} \appear{\Alert[fuchsia]{explains why real materials break and transform} so much more easily} \appear{than what is predicted by a simple (but reasonable) theory based on ideal crystalline solids}

\end{itemize}

\endofslide 
%\endofsection
\end{frame}
%%%%%%%%%%%%%%%

%%%%%%%%%%%%%%%
\begin{frame}
\frametitle{\texorpdfstring{\culine[\sectiontitleulinecolor]{\sectiontitle}}{\sectiontitle} }
\framesubtitle{Crystal defects}
%Crystal defects
\begin{itemize}
	
\item Dislocation moves to the right as the atoms in the layer under it successively shift their bonds with those of the upper layer\appear{, \Alert[fuchsia]{one line at the time}} \appear{(see Figures a–c)}
\appear{\xdef\loc{south}%
\begin{tikzpicture}[remember picture,overlay]% anchor=south
    \node (A) [yshift=-0.05*\figYshift,xshift=-0*\figXshift, anchor=\loc] at (current page.\loc) {\includegraphics[page=1,width=14cm]{Figures_booklet/SOM_10_Bonding_figures/Tikz_SOM_Bonding_Mechanical_dislocation_movement_fixed}};%Tikz_SOM_Bonding_Mechanical_dislocation_movement
\end{tikzpicture}}
\item Process requires dramatically less energy for a given instant of time\appear{, than what would be required to move several (thousands) of lines at a time}

\item These kinds of \CDAlert[fuchsia]{(plastic) deformations} can be controlled\appear{, by intentionally adding defects into the material (Figure d) }
	
\end{itemize}



\endofslide 
%\endofsection
\end{frame}
%%%%%%%%%%%%%%%

%%%%%%%%%%%%%%%
\begin{frame}
\frametitle{\texorpdfstring{\culine[\sectiontitleulinecolor]{\sectiontitle}}{\sectiontitle} }
\framesubtitle{Crystal lattice and the lattice points }

\begin{itemize}[itemsep=1.6mm]
	\item Crystalline solids are \Alert[red]{categorized according to their symmetries into} 7 crystal systems and \Alert[red]{14 Bravais lattices} 
	
	\item The structure of the crystal\appear{, i.\,e. \CDAlert[fuchsia]{crystal lattice}}\appear{, can be created by repeating the unit cell }
	
	\item The unit cell is given by \CDAlert[blue]{three translation vectors $\mathbf{a}_i$}\appear{, where $i=1$, 2, 3}\appear{, so that if one looks at the lattice from two different but equivalent points $\mathbf{r}^{\prime}$} \appear{and $\mathbf{r}$}\appear{, then these \CDAlert[fuchsia]{lattice points} are related to each other via:}
\[ 
 \mathbf{r}^{\prime}  \equal \appear{\mathbf{r}}  \plus \appear{m_{1} \mathbf{a}_1} \plus \appear{m_{2} \mathbf{a}_2}  \plus \appear{m_{3} \mathbf{a}_3} \appear{\,, \qquad where~m_{i}}\appear{=integers} 
 \]

\item The choice for the unit cell is not unique\appear{, it can be chosen by many different ways}

\item \Alert[red]{The smallest cell is the \CDAlert[red]{primitive cell}}\appear{, and the lengths of its edges are called \CDAlert[tunired]{lattice constants}}
%\item placed in each lattice point in same direction and posture
\end{itemize}

\endofslide 
%\endofsection
\end{frame}
%%%%%%%%%%%%%%%

%%%%%%%%%%%%%%%
\begin{frame}
\frametitle{\texorpdfstring{\culine[\sectiontitleulinecolor]{\sectiontitle}}{\sectiontitle} }
\framesubtitle{Crystal structures, cubic structures}

%\begin{itemize}
%	\item \CDAlert[red]{Body-centered cubic (bcc)}:
%\item[] \begin{itemize}
%	\item A repetition of cubes with atoms placed at the corners and at the center of the cube
%\item Solids from the periodic table's first column: Li, Na, K, Rb, and Cs
%\item Transition metals, such as: Fe, Cr, W
%\end{itemize}
%\item \CDAlert[blue]{Face-centered cubic (fcc)}:
%\item[] \begin{itemize}
%	\item A repetitive cubic structure, with atoms placed both at the corners of the cube and at the centers of the cubic faces
%	\item Noble gases (except helium) 
%	\item Transition metals such as: Cu, Ag, Au, Ni, Pb
%\end{itemize}
%\end{itemize}

\vspace{1mm}
\begin{minipage}{10.5cm}
\begin{itemize}%[itemsep=1.6mm] 
\item \CDAlert[red]{Body-centered cubic (bcc)}:
\begin{itemize}
	\item A repetition of cubes with atoms placed at the corners \appear{and at the 			center of the cube}
	
	\item Solids from the periodic table's first column: \appear{Li}\appear{, Na}\appear{, K}\appear{, Rb, and Cs}
	\item Transition metals\appear{, such as: Fe, Cr, W}
	\end{itemize}

\item \CDAlert[blue]{Face-centered cubic (fcc)}:
\begin{itemize}
	\item A repetitive cubic structure\appear{, with atoms placed both at the corners of the cube} \appear{and at the centers of the cubic facets}
	\item Noble gases \appear{(except helium)}
	\item Transition metals \appear{such as: Cu, Ag, Au, Ni, Pb}
	\end{itemize}
\end{itemize}
\end{minipage}
%% 
\xdef\loc{south east}
\begin{tikzpicture}[remember picture,overlay] % anchor=south
    \node (A) [yshift=-0.1*\figYshift,xshift=\figXshift, anchor=\loc] at (current page.\loc) {\includegraphics[page=1,height=9cm]{Figures_booklet/SOM_10_Bonding_figures/SOM_Bonding_Figure_22.png}}; %scale=\figscale. % 8cm max height
\end{tikzpicture}

\endofslide 
%\endofsection
\end{frame}
%%%%%%%%%%%%%%%

%%%%%%%%%%%%%%%
\begin{frame}
\frametitle{\texorpdfstring{\culine[\sectiontitleulinecolor]{\sectiontitle}}{\sectiontitle} }
\framesubtitle{Crystal structures, hexagonal closest packing}

\vspace{1mm}
\begin{minipage}{10.5cm}
\begin{itemize}%[itemsep=1.6mm] 
\item \CDAlert[red]{Hexagonal closest packing (hcp)}:
\begin{itemize}
	\item Six-fold rotational symmetry
\item Elements throughout the periodic table:
\item He, Be, Cd, Ce, Mg, Os, Zn, Zr, and many rare earths

\end{itemize}

\item \CDAlert[red]{Simple cubic (sc)}:
\begin{itemize}
	\item A repetitive cubic structure, with atoms placed at the corners of the cube
\end{itemize}

\vspace{3mm}

\item The above examples assume: \appear{\CDAlert[fuchsia]{lattice point = atom}}

\item In real life, several atoms/molecules can occupy each lattice point of the crystalline structure\appear{, resulting in 3D structures that are \CDAlert[blue]{non-trivial to visualize}}
\end{itemize}
\end{minipage}

%% 
\xdef\loc{south east}
\begin{tikzpicture}[remember picture,overlay] % anchor=south
    \node (A) [yshift=-0.1*\figYshift,xshift=\figXshift, anchor=\loc] at (current page.\loc) {\includegraphics[page=1,height=9cm]{Figures_booklet/SOM_10_Bonding_figures/SOM_Bonding_Figure_22.png}}; %scale=\figscale. % 8cm max height
\end{tikzpicture}


\endofslide 
%\endofsection
\end{frame}
%%%%%%%%%%%%%%%

%%%%%%%%%%%%%%%
\begin{frame}
\frametitle{\texorpdfstring{\culine[\sectiontitleulinecolor]{\sectiontitle}}{\sectiontitle} }
\framesubtitle{3D point groups}

\begin{itemize}
	\item In summary, 7 crystal systems and 14 Bravais lattices of crystalline solids (\Alert[red]{link}{https://en.wikipedia.org/wiki/Space_group\#Table_of_space_groups_in_3_dimensions}):
	
%	\item Maybe only show the 2D groups here and mention that 3Ds will be covered in the future courses...?
\end{itemize}

%https://en.wikipedia.org/wiki/Space_group#Table_of_space_groups_in_3_dimensions
\xdef\loc{south}
\begin{tikzpicture}[remember picture,overlay] % anchor=south
    \node (A) [yshift=-0.1*\figYshift,xshift=0*\figXshift, anchor=\loc] at (current page.\loc) {\includegraphics[page=1,height=8.2cm]{Figures_booklet/SOM_10_Bonding_figures/Tikz_SOM_Bonding_Lattices_fixed}}; %scale=\figscale. % 8cm max height
\end{tikzpicture}

\endofslide 
%\endofsection
\end{frame}
%%%%%%%%%%%%%%%

%%%%%%%%%%%%%%%
\begin{frame}
\frametitle{\texorpdfstring{\culine[\sectiontitleulinecolor]{\sectiontitle}}{\sectiontitle} }
\framesubtitle{Covalent solids}

\begin{itemize}
	\item \CDAlert[red]{Covalently bonded solid} has bonds similar to covalent bonds of molecules described previously\appear{, i.\,e. electrons are shared between the atoms}

\item In covalent solids each atom shares covalent bonds with other surrounding atoms 
 \implies{ unbroken network of strong bonds}
 
 \item \Alert[red]{Such solids are hard and have high melting points}\appear{, they are poor electrical conductors} \appear{(\CDAlert[blue]{dielectrics} or \CDAlert[blue]{semiconductors})} \appear{because all the \CDAlert[tuniblue]{valence electrons are locked into bonds} between adjacent atoms}

\end{itemize}

\endofslide 
%\endofsection
\end{frame}
%%%%%%%%%%%%%%%


%%%%%%%%%%%%%%%
\begin{frame}
\frametitle{\texorpdfstring{\culine[\sectiontitleulinecolor]{\sectiontitle}}{\sectiontitle} }
\framesubtitle{Covalent solids}

\begin{itemize}
	
 \item \CDAlert[violet]{Diamond} is an exemplary covalent solid\appear{, and forms in $fcc$} \appear{but has two carbon atoms in each lattice point} \implies{ Structure is special and known as the \CDAlert[violet]{diamond structure} }

\item In other words\appear{, diamond consists of two inter-penetrating fcc structures} 

\item Tight binding also results in the highest atomic cohesive energy of all solids, $7.37 \mathrm{~eV}$ per atom

\end{itemize}

\vspace{2.5mm}
\begin{minipage}{8cm}
\begin{itemize}%[itemsep=1.6mm] 
	\item Note that carbon has several other well defined crystalline structures\appear{, such as \CDAlert[tunired]{graphite}}\appear{, solid \CDAlert[fuchsia]{fullerenes}} \appear{and \CDAlert[red]{graphene}}
\end{itemize}
\end{minipage}

\xdef\loc{south east}
\begin{tikzpicture}[remember picture,overlay] % anchor=south
    \node (A) [yshift=-0.1*\figYshift,xshift=\figXshift, anchor=\loc] at (current page.\loc) {\includegraphics[page=1,height=5cm]{Figures_booklet/SOM_10_Bonding_figures/Tikz_SOM_Bonding_diamond_fixed}}; %scale=\figscale. % 8cm max height
\end{tikzpicture}

\endofslide 
%\endofsection
\end{frame}
%%%%%%%%%%%%%%%


%%%%%%%%%%%%%%%
\begin{frame}
\frametitle{\texorpdfstring{\culine[\sectiontitleulinecolor]{\sectiontitle}}{\sectiontitle} }
\framesubtitle{Some of the many allotropes of carbon}


\xdef\loc{south west}
\begin{tikzpicture}[remember picture,overlay] % anchor=south
    \node (A) [yshift=-0.1*\figYshift,xshift=-1*\figXshift, anchor=\loc] at (current page.\loc) {\includegraphics[page=1,width=5cm]{Figures_booklet/SOM_10_Bonding_figures/Tikz_SOM_Bonding_figures_graphene_graphite}}; %scale=\figscale. % 8cm max height
\end{tikzpicture}

\xdef\loc{west}
\begin{tikzpicture}[remember picture,overlay] % anchor=south
    \node (A) [yshift=-1.5*\figYshift,xshift=-1*\figXshift, anchor=\loc] at (current page.\loc) {\includegraphics[page=2,width=5.5cm]{Figures_booklet/SOM_10_Bonding_figures/Tikz_SOM_Bonding_figures_graphene_graphite}}; %scale=\figscale. % 8cm max height
\end{tikzpicture}
%Tikz_SOM_Bonding_figures_graphene_graphite

\xdef\loc{south}
\begin{tikzpicture}[remember picture,overlay] % anchor=south
    \node (A) [yshift=-0.1*\figYshift,xshift=0*\figXshift, anchor=\loc] at (current page.\loc) {\includegraphics[page=1,height=8.5cm]{Figures_booklet/SOM_10_Bonding_figures/SOM_Bonding_Figure_58.png}}; %scale=\figscale. % 8cm max height
\end{tikzpicture}

\xdef\loc{south east}
\begin{tikzpicture}[remember picture,overlay] % anchor=south
    \node (A) [yshift=-0.1*\figYshift,xshift=\figXshift, anchor=\loc] at (current page.\loc) {\includegraphics[page=1,height=6cm]{Figures_booklet/SOM_10_Bonding_figures/SOM_Bonding_Figure_57.png}}; %scale=\figscale. % 8cm max height
\end{tikzpicture}

\endofslide 
%\endofsection
\end{frame}
%%%%%%%%%%%%%%%

%%%%%%%%%%%%%%%
\begin{frame}
\frametitle{\texorpdfstring{\culine[\sectiontitleulinecolor]{\sectiontitle}}{\sectiontitle} }
\framesubtitle{Ionic solids}

\begin{itemize}
	\item In the regular table salt $\mathrm{NaCl}$\appear{, the ions Na$^+$ and Cl$^-$ are both spherically symmetric}\appear{, but Cl$^-$ ion is approximately twice as large as Na$^+$}

\item The minimum of the potential energy is obtained when an ion of either kind has six \CDAlert[blue]{nearest neighbours} of the other kind 
\implies{ $\mathrm{NaCl}$ has a face-centered cubic (fcc) structure, with \CDAlert[red]{each lattice} \CDAlert[red]{point consisting of 1 Na and 1 Cl atom}} 
\end{itemize}

\vspace{2.5mm}
\begin{minipage}{8cm}
\begin{itemize}%[itemsep=1.6mm] 
	\item In other words, the fcc structure has, in fact, two fcc lattices within each other
\item Incidentally, table salt (NaCl), is quite important solid: \CDAlert[tunired]{every living} \CDAlert[tunired]{creature on Earth needs NaCl}

\end{itemize}
\end{minipage}

\xdef\loc{south east}
\begin{tikzpicture}[remember picture,overlay] % anchor=south
    \node (A) [yshift=-0.25*\figYshift,xshift=\figXshift, anchor=\loc] at (current page.\loc) {\includegraphics[page=1,height=5cm]{Figures_booklet/SOM_10_Bonding_figures/Tikz_SOM_Bonding_NaCl_fixed}}; %scale=\figscale. % 8cm max height
\end{tikzpicture}

\endofslide 
%\endofsection
\end{frame}
%%%%%%%%%%%%%%%

%%%%%%%%%%%%%%%
\begin{frame}
\frametitle{\texorpdfstring{\culine[\sectiontitleulinecolor]{\sectiontitle}}{\sectiontitle} }
\framesubtitle{Ionic solids and Madelung constant}

\begin{itemize}
	\item Note, that $\mathrm{Na}$ and $\mathrm{Cl}$ are not paired into $\mathrm{NaCl}$ in the solid structure 
	
	\item But the electric energy for the ion pair can still be written:
\[ 
 U_{\text {attraction}}(r)  \equal \appear{-\alpha} \frac{1}{4 \pi \varepsilon_{0}} \frac{e^{2}}{r}  \equal \appear{-\alpha} \frac{k e^{2}}{r} \appear{\,,\qquad \text{''} \Frac{1}{4 \pi \varepsilon_{0}} \Equiv k \text{''} 
} \]%\text{ "}
\appear{Where $r$ is the distance between the neighbouring ions} \appear{$(0.282 \mathrm{~nm}$ for $\mathrm{NaCl}$)}

\item $\alpha$ is so-called \CDAlert[red]{Madelung constant}\appear{, which takes into account the electrostatic repulsion}\appear{/attraction of the neighbours} 

\item If only the nearest 6 neighbours with opposite charge had an effect\appear{, then the value of $\alpha$ would be 6} 

\end{itemize}

\endofslide 
%\endofsection
\end{frame}
%%%%%%%%%%%%%%%


%%%%%%%%%%%%%%%
\begin{frame}
\frametitle{\texorpdfstring{\culine[\sectiontitleulinecolor]{\sectiontitle}}{\sectiontitle} }
\framesubtitle{Madelung constant}

\begin{itemize}
	
\item But in the lattice structure there are also 12 ions with charge of the same sign at distance $2^{1 / 2} r$\appear{, and 8 ions of the opposite charge at $3^{1 / 2} r$ etc.} \appear{Thus, the total contribution is:}
\[ 
 \alpha  \equal \appear{6-} \frac{12}{\sqrt{2}} \plus \frac{8}{\sqrt{3}}  \minus \frac{6}{2} \plus \frac{20}{\sqrt{5}}  \minus \appear{\ldots}
 \]
\appear{resulting in a sum that does not converge!} 

\item But if we group the terms in a clever way\appear{, it is possible to make it converge}\appear{, although slowly, and the result is $\alpha=1.7476$}
\end{itemize}

\endofslide 
%\endofsection
\end{frame}
%%%%%%%%%%%%%%%

%%%%%%%%%%%%%%%
\begin{frame}
\frametitle{\texorpdfstring{\culine[\sectiontitleulinecolor]{\sectiontitle}}{\sectiontitle} }
\framesubtitle{Madelung constant}

\begin{itemize}
	\item The geometric details of some common solids are given in the table
\end{itemize}
\appear{\xdef\loc{south east}%
\begin{tikzpicture}[remember picture,overlay] % anchor=south
    \node (A) [yshift=-0.1*\figYshift,xshift=\figXshift, anchor=\loc] at (current page.\loc) {\includegraphics[page=1,height=8.2cm]{Figures_booklet/SOM_Tipler_Solid_state_physics_figures/SOM_Tipler_SSP_figure_5.png}};%scale=\figscale. % 8cm max height
\end{tikzpicture}}

\vspace{2.5mm}
\begin{minipage}{5.5cm}
\begin{itemize}%[itemsep=1.6mm] 
\item Note, that Madelung constant is applied only for ionic solids 

\item And, for other (ionic) crystal structures\appear{, the value for $\alpha$ will be different}
\end{itemize}
\end{minipage}


\endofslide 
%\endofsection
\end{frame}
%%%%%%%%%%%%%%%

%%%%%%%%%%%%%%%
\begin{frame}
\frametitle{\texorpdfstring{\culine[\sectiontitleulinecolor]{\sectiontitle}}{\sectiontitle} }
\framesubtitle{Ionic solids}

\begin{itemize}[itemsep=1.8mm]
	\item When the ions are close together\appear{, the wave functions of the electrons start overlapping} \appear{and an repulsive force appears due to the \CDAlert[red]{Pauli exclusion principle}} 
	
	\item In this case the potential energy will be given by
\[ 
 U(r)  \equal \appear{-\alpha} \frac{k e^{2}}{r}  \plus \frac{A}{r^{n}} 
 \]

\item And by differentiating the energy to obtain a force balance at the equilibrium separation $r=r_{0}$\appear{, we can solve the constant $A$ to be}
\[ 
 A  \equal  \frac{\alpha k e^{2} r_{0}^{n-1}}{n} 
 \]

\item So the potential energy can be written as
\[ 
 U(r)  \equal \appear{-\alpha} \frac{k e^{2}}{r_{0}} \appear{\textrm{[[} \Frac{r_{0}}{r}} \appear{-\Frac{1}{n}} \appear{\textrm{((}} \frac{r_{0}}{r} \appear{\textrm{))}^{n}} \appear{\textrm{]]} }
 \]

\end{itemize}

\endofslide 
%\endofsection
\end{frame}
%%%%%%%%%%%%%%%

%%%%%%%%%%%%%%%
\begin{frame}
\frametitle{\texorpdfstring{\culine[\sectiontitleulinecolor]{\sectiontitle}}{\sectiontitle} }
\framesubtitle{Ionic solids}

\begin{itemize}[itemsep=1.8mm]

\item Which, for the equilibrium separation $r=r_{0}$ will result in\vspace{-1.5mm}

\[ 
 U \textrm{((}r_{0} \textrm{))}  \equal \appear{-\alpha} \frac{k e^{2}}{r_{0}} \appear{\textrm{[[}1-\Frac{1}{n}\textrm{]]} }
 \]

	\item The equilibrium separation can be determined experimentally using \CDAlert[red]{x-ray diffraction} \appear{(or by \CDAlert[blue]{computing it from the crystal density}) }%
	\vspace{-2.5mm}
	
	\implies{ Dissociation energy or lattice energy of the ionic crystal gives $n$ }
	\item Dissociation energy of $\mathrm{NaCl}$ is $770 \mathrm{\;kJ} / \mathrm{mol}\appear{=7.98$\;eV/\,ion pair}\appear{, resulting in $n=9.35 \approx 9$}
\end{itemize}

\vspace{1.6mm}
\begin{minipage}{8.8cm}
\begin{itemize}[itemsep=1.6mm]
	 
\item[] \begin{itemize}
	\item Here, the value $7.98$\;eV is now the energy needed to remove a Na$^+$ and Cl$^-$ ion pair from the crystal 
	\item Electron affinity of $\mathrm{Cl}$ is $3.62$\;eV
	\item Forming Na from Na$^+$ \appear{releases $5.15 \mathrm{\;eV}$} 
\end{itemize}

\end{itemize}
\end{minipage}

\xdef\loc{south east}
\begin{tikzpicture}[remember picture,overlay] % anchor=south
    \node (A) [yshift=-0.1*\figYshift,xshift=\figXshift, anchor=\loc] at (current page.\loc) {\includegraphics[page=1,height=3.85cm]{Figures_booklet/SOM_10_Bonding_figures/SOM_Bonding_Figure_16.png}}; %scale=\figscale. % 8cm max height
\end{tikzpicture}

\endofslide 
%\endofsection
\end{frame}
%%%%%%%%%%%%%%%

%%%%%%%%%%%%%%%
\begin{frame}
\frametitle{\texorpdfstring{\culine[\sectiontitleulinecolor]{\sectiontitle}}{\sectiontitle} }
\framesubtitle{Ionic solids}

\begin{itemize}

\item Putting things together, the energy needed to remove neutral $\mathrm{Na}$ \appear{and $\mathrm{Cl}$ pair of atoms} \appear{from the crystal will be} $\appear{7.98 \mathrm{\;eV}} \plus \appear{3.62 \mathrm{\;eV}}  \minus \appear{5.15 \;\mathrm{eV}} \equal \appear{6.45 \mathrm{\;eV}}$

\item So, eventually the so-called \CDAlert[red]{cohesive energy of a crystal}\appear{, which is reported for atoms or per atomic pairs}\appear{, will be $\Frac{6.45 \mathrm{\;eV}}{2}=3.225$~eV/atom}\appear{, which value is close to literature value of $3.19 \mathrm{\;eV}/$atom} 

\item \CDAlert[tunired]{A large cohesive energy implies high melting point}

	\item In addition to $\mathrm{NaCl}$\appear{, most ionic crystals have a face-centered cubic structure $(f c c)$:} \appear{LiF, KF, KCl, KI, AgCl, ...}

\item Another ionic crystal structure is simple cubic (sc)\appear{, such as in the case of CsCl}\appear{, TlI, TlBr, LiHg, NH$_{4}$,  Cl, ...}

\item Madelung constant for an sc structure is $\alpha=1.7627$ 
\end{itemize}


\endofslide 
%\endofsection
\end{frame}
%%%%%%%%%%%%%%%

%%%%%%%%%%%%%%%
\begin{frame}
\frametitle{\texorpdfstring{\culine[\sectiontitleulinecolor]{\sectiontitle}}{\sectiontitle} }
\framesubtitle{Ionic solids and Madelung constant example}

\begin{itemize}
	\item \CDAlert[red]{Example}: Calculate the Madelung constant for the two-dimensional crystal shown in figure%
	\appear{\xdef\loc{south east}%
\begin{tikzpicture}[remember picture,overlay]% anchor=south
    \node (A) [yshift=-0.15*\figYshift,xshift=\figXshift, anchor=\loc] at (current page.\loc) {\includegraphics[page=1,height=6cm]{Figures_booklet/SOM_Tipler_Solid_state_physics_figures/SOM_Tipler_SSP_figure_9.png}};%scale=\figscale. % 8cm max height
\end{tikzpicture}}
 [Hint: calculate only first 4 terms]
	
\end{itemize}

\vspace{1.6mm}
\begin{minipage}{6.5cm}
\begin{itemize}[itemsep=1.6mm]
	 
\item[] \begin{itemize}
	\item[] $$
\begin{aligned}
&\appear{U_{\text {attraction }}} \appear{=}\appear{-k e^{2}} \appear{\textrm{((}}\frac{4}{r} \minus \frac{4}{\sqrt{2} r} \minus \frac{4}{2 r} \plus \frac{8}{\sqrt{5} r} \appear{\textrm{))}} \\
&\appear{U_{\text {attraction }}}  \equal  \minus \frac{k e^{2}}{r} \appear{\textrm{((}4} \minus \frac{4}{\sqrt{2}} \minus \appear{2} \plus \frac{8}{\sqrt{5}} \appear{\textrm{))}}
\end{aligned}
$$
\item[] Answer: \appear{$\alpha \approx 2.749$} 
\end{itemize}

\end{itemize}
\end{minipage}

\endofslide 
%\endofsection
\end{frame}
%%%%%%%%%%%%%%%

%%%%%%%%%%%%%%%
\begin{frame}
\frametitle{\texorpdfstring{\culine[\sectiontitleulinecolor]{\sectiontitle}}{\sectiontitle} }
\framesubtitle{Metallic solids}

\begin{itemize}[itemsep=1.7mm]
	\item Except for noble gases, all elements have valence electrons
	
	\item But in most cases no arrangement of covalent bonds can join all of a given atom's valence electrons to those of the atoms surrounding it 
	
	\item An element or compound \CDAlert[tunired]{with ''leftover'' valence electrons} \appear{forms a \CDAlert[red]{metallic solid}}

\item A virtual continuum of states is produced\appear{, whose energies are lower than those of isolated atoms }
\end{itemize}

\vspace{1.6mm}
\begin{minipage}{7.5cm}
\begin{itemize}[itemsep=1.6mm]
	 \item Shared by all the atoms\appear{, the valence electrons move rather freely in the crystal lattice}\appear{, forming an \CDAlert[red]{electron gas}} 
\implies{ they become \CDAlert[tunired]{conduction electrons} }

\item \Alert[red]{Metallic solids are \CDAlert[red]{good conductors} of \CDAlert[red]{electricity and heat}}

\end{itemize}
\end{minipage}


\xdef\loc{south east}
\begin{tikzpicture}[remember picture,overlay] % anchor=south
    \node (A) [yshift=-0.15*\figYshift,xshift=\figXshift, anchor=\loc] at (current page.\loc) {\includegraphics[page=1,height=4.4cm]{Figures_booklet/SOM_10_Bonding_figures/Tikz_SOM_Bonding_e_conduction_schematic_fixed}}; %scale=\figscale. % 8cm max height
\end{tikzpicture}

\endofslide 
%\endofsection
\end{frame}
%%%%%%%%%%%%%%%

%%%%%%%%%%%%%%%
\begin{frame}
\frametitle{\texorpdfstring{\culine[\sectiontitleulinecolor]{\sectiontitle}}{\sectiontitle} }
\framesubtitle{Metallic solids}

\begin{itemize}
	\item Let's consider solid lithium as an example

\item[] \begin{itemize}
	\item Lithium atom has an electron structure of $\appear{1 s^{2}} \appear{2 s}$ 

\item The radial wave function of the $2 s$ electron "sees" a hydrogen-like core \appear{consisting of the nucleus and the completed 1s shell}\appear{, and is of form}
\[ 
 \psi_{2,0} \propto \appear{R_{n=2, \ell=0}} \qquad \appear{\&} \qquad \appear{\psi_{2,0}} \equal \appear{C_{2,0}} \appear{\textrm{((}2} \minus \frac{r}{a_{0}} \appear{\textrm{))} }\appear{e^{-r / 2 a_{0}} }
 \]
\appear{where the $C_{2,0}$ is a normalization constant} \appear{$ \textrm{((}C_{2,0}= \textrm{((}2 a_{0} \textrm{))}^{-3 / 2} \textrm{))}$}\appear{, and $a_{0}=$~Bohr radius}

\end{itemize}

\end{itemize}

\endofslide 
%\endofsection
\end{frame}
%%%%%%%%%%%%%%%

%%%%%%%%%%%%%%%
\begin{frame}
\frametitle{\texorpdfstring{\culine[\sectiontitleulinecolor]{\sectiontitle}}{\sectiontitle} }
\framesubtitle{Metallic solids}

\begin{itemize}
	\item Let's consider solid lithium as an example:

\item[] \begin{itemize}
	\item The probability densities are plotted in the figure:
	\item[] \begin{itemize}
			\item[(a)] is for an isolated lithium atom
			\item[(b)] is for the metallic solid
		\end{itemize}
\end{itemize}
\end{itemize}

\vspace{1.6mm}
\begin{minipage}{4.5cm}
\begin{itemize}[itemsep=1.6mm]
\item[] \begin{itemize}
	 \item Potential energy of the electrons has been reduced\appear{, and spatial locations of the electrons are more peaked around $\pm 0.3 \mathrm{\;nm}$}

\end{itemize}

	
\end{itemize}
\end{minipage}

\xdef\loc{south east}
\begin{tikzpicture}[remember picture,overlay] % anchor=south
    \node (A) [yshift=-0.15*\figYshift,xshift=\figXshift, anchor=\loc] at (current page.\loc) {\includegraphics[page=1,height=6cm]{Figures_booklet/SOM_Tipler_Solid_state_physics_figures/SOM_Tipler_SSP_figure_11.png}}; %scale=\figscale. % 8cm max height
\end{tikzpicture}

\endofslide 
%\endofsection
\end{frame}
%%%%%%%%%%%%%%%

%%%%%%%%%%%%%%%
\begin{frame}
\frametitle{\texorpdfstring{\culine[\sectiontitleulinecolor]{\sectiontitle}}{\sectiontitle} }
\framesubtitle{Molecular solids}

\begin{itemize}
	\item A molecular solid consists of discrete molecules 
	
	\item The cohesive forces binding the molecules together are \CDAlert[red]{Van der Waals forces:}\appear{ \CDAlert[tunired]{dipole–dipole interaction}}\appear{, \CDAlert[red]{London force} etc. }
	
	\item Even though electrons sit on the covalent molecules so tightly that there is no incentive to share the electrons\appear{, there are still attractive forces}\appear{, when the molecules are brought close to each other}

\item For symmetric covalent molecules \appear{(e.g. $\mathrm{N}_{2}$ or $\mathrm{H}_{2}$)} \appear{with no permanent electric dipole}\appear{, the molecules tend to be polarised} 
\implies{ The \CDAlert[red]{induced dipole–dipole attraction} is known as \CDAlert[red]{London force} }

\end{itemize}

\endofslide 
%\endofsection
\end{frame}
%%%%%%%%%%%%%%%

%%%%%%%%%%%%%%%
\begin{frame}
\frametitle{\texorpdfstring{\culine[\sectiontitleulinecolor]{\sectiontitle}}{\sectiontitle} }
\framesubtitle{Molecular solids}

\begin{itemize}
	\item If the molecules have permanent dipole moments \appear{(e.g. $\mathrm{NH}_{3}$ and $\mathrm{H}_{2} \mathrm{O}$)}\appear{, an additional, somewhat stronger}\appear{, \CDAlert[red]{dipole–dipole attraction} becomes important} 
	\implies{ This is exceptionally strong for \CDAlert[blue]{water}, where the attraction\\
	 is called the \CDAlert[blue]{hydrogen bond} }

\item \CDAlert[red]{Molecular solids}:
\item[] \begin{itemize}
	\item Are relatively soft
\item Have low melting points \appear{$ \textrm{((}\!<200^{\circ}C \textrm{))}$}
\item Have poor conductivity of heat \appear{and electricity}
\item \CDAlert[blue]{Water ice $H_2 O$} \appear{and dry ice $CO_2$} \appear{are representative examples}
\end{itemize}

\end{itemize}

\endofslide 
%\endofsection
\end{frame}
%%%%%%%%%%%%%%%

%%%%%%%%%%%%%%%
\begin{frame}
\frametitle{\texorpdfstring{\culine[\sectiontitleulinecolor]{\sectiontitle}}{\sectiontitle} }
\framesubtitle{Summary}

\footnote{Encyclopaedia Britannica: metallic bond (chemistry)}
% https://www.britannica.com/science/metallic-bond
\xdef\loc{center}
\begin{tikzpicture}[remember picture,overlay] % anchor=south
    \node (A) [yshift=0.35*\figYshift,xshift=0*\figXshift, anchor=\loc] at (current page.\loc) {\includegraphics[page=1,width=14cm]{Figures_booklet/SOM_bonding_types_of_bonds.png}}; %scale=\figscale. % 8cm max height
\end{tikzpicture}

%\endofslide 
%\endofsection
\end{frame}
%%%%%%%%%%%%%%%